\section{Глава~\ref{sec:dendrites}. Число дендритов}

Строение классов и числа входов измерены анатомически и электрическими записями; оценка~\eqref{eq:dendriteload} --- простая физика конденсатора, а вычислительное объяснение числа входов опирается на теорию и модели.

{\footnotesize
\setlength{\tabcolsep}{3pt}\renewcommand{\arraystretch}{1.15}
\begin{longtable}{>{\raggedright}p{0.5cm}>{\raggedright}p{4.0cm}>{\raggedright}p{5.6cm}>{\raggedright}p{2.3cm}>{\raggedright\arraybackslash}p{1.7cm}}
\toprule
№ & Утверждение & Опыт и что измерено & Источник & Статус \\
\midrule\endfirsthead
\toprule
№ & Утверждение & Опыт и что измерено & Источник & Статус \\
\midrule\endhead
1 & Удельная ёмкость мембраны $c_m\approx1$~мкФ/см$^2$ & С тела нейрона снимали мембрану на пипетку, подавали ступеньку напряжения и по спаду ёмкостного тока находили полную ёмкость, затем делили на площадь: $0.9$~мкФ/см$^2$ у трёх классов нейронов & \cite{Gentet2000} & измерено \\
2 & Время заряда $t\approx c_mA\Delta V/I$~\eqref{eq:dendriteload}: большому нейрону нужны десятки одновременных входов, маленькому хватает одного крупного синапса & Оценка по закону заряда конденсатора; числа ($20$ и $300$~пФ, $0.1$ и несколько нА) --- порядки величин & вывод в главе & теория \\
3 & Один огромный синапс даёт ток в несколько наноампер и сразу доводит маленький нейрон до порога & Одновременная запись окончания и получателя в срезе: ток одного синапса $2$--$13$~нА, каждый импульс входа вызывает надпороговый ответ, задержка около $0.4$~мс при $36^\circ$C; получатель выдаёт \nt{ГЛИЦ} & \cite{Borst1995,BorstSoria2012} & измерено \\
4 & Ноль дендритов: у чувствительных нейронов осязания и боли импульс идёт от датчика в спинной мозг мимо тела клетки & Строение (аксон, делящийся надвое у тела) установлено анатомически. То, что проведение не требует возбудимости тела, показано на проверенной модели такого нейрона: без возбудимого тела импульс проходит развилку так же надёжно & \cite{AmirDevor2003} & теория \\
5 & Один дендрит: обонятельный нейрон несёт один тип рецептора и передаёт одну входную линию & Обзор опытов с генами рецепторов: каждый обонятельный нейрон выражает один ген рецептора из большого семейства & \cite{Mombaerts2004} & измерено \\
6 & Два дендрита: у детекторов направления на звук входы от левого и правого уха садятся на разные дендриты & Анатомия и записи у млекопитающих, собранные в обзоре: входы от двух ушей разнесены по двум противоположным дендритам & \cite{Grothe2010} & измерено \\
7 & Разнесение входов по двум дендритам делает «И» строже & Показано на модели: насыщение~\eqref{eq:shunt} на одном дендрите не даёт входам с одной стороны дойти до порога, и детекция совпадений улучшается. Прямого опыта, где бы меняли раскладку входов, нет & \cite{AgmonSnir1998} & теория \\
8 & Около $7\cdot10^{10}$ маленьких \nt{ГЛУТ}-нейронов цепи обучения движениям, по $\sim4$ входа у каждого & Подсчёт клеток изотропным фракционированием: в этой части мозга человека около $69$ млрд нейронов. Запись одной такой клетки у живой крысы: на касание приходят пачки дискретных токов от немногих входов, и клетка срабатывает, когда они складываются & \cite{Azevedo2009,Chadderton2004} & измерено \\
9 & Пороговый элемент с $n$ входами разделяет не больше $P_{\max}\approx2n$ случайных образов~\eqref{eq:cover} & Классический результат комбинаторной геометрии. Оговорка: при одних лишь положительных весах и с запасом на шум ёмкость меньше; для выходного нейрона цепи обучения движениям оценено $\sim4\cdot10^4$ соответствий, а не $2\cdot10^5$ & \cite{Cover1965,Brunel2004} & теория \\
10 & Смесители с малым числом входов нужны, чтобы расширить вход; «четыре» близко к оптимуму & Теория: размерность представления, которую даёт слой смесителей, максимальна примерно при том числе входов, которое наблюдается анатомически; исходная гипотеза расширения --- в классических работах & \cite{Marr1969,Albus1971,LitwinKumar2017} & гипотеза \\
11 & Большие деревья: $10^4$--$10^5$ входов & Подсчёт синапсов по электронным снимкам: $\approx1.7\cdot10^5$ синапсов на одном выходном \nt{ГАМК}-нейроне цепи обучения движениям у крысы; около $30\,000$ возбуждающих и $1\,700$ тормозящих входов на \nt{ГЛУТ}-нейроне гиппокампа & \cite{NapperHarvey1988,Megias2001} & измерено \\
12 & Ветви большого дерева --- отдельные подсхемы со своими порогами; нейрон --- маленькая многослойная сеть & Точечная стимуляция двух мест на тонких дендритах: входы на одной ветви складываются по сигмоиде, на разных ветвях --- линейно. В дендритах нейронов коры человека найдены градуальные кальциевые импульсы, позволяющие одной клетке решать линейно неразделимую задачу. Модель: чтобы воспроизвести один нейрон, нужна глубокая сеть & \cite{Polsky2004,Gidon2020,Beniaguev2021} & измерено \\
13 & Дендриты-выходы: каждая ветвь \nt{ГАМК}-нейрона сетчатки --- отдельный детектор направления & Двухфотонная запись кальция в отдельных ветвях: сигнал в ветви избирателен к движению от тела клетки к кончику, а напряжение тела --- нет & \cite{Euler2002} & измерено \\
14 & Число входов следует за задачей (утверждение~\ref{prop:dendrites}) & Строение классов измерено (строки 3--13); что число входов выбрано ради скорости или классификации, показано лишь теорией и моделями для отдельных классов & \cite{LitwinKumar2017,AgmonSnir1998} & гипотеза \\
\bottomrule
\end{longtable}}
